% ====================================================================
% 부록 B — 문건→코드 요구명세 (doc-leads: 코드가 이 문건을 구현할 때 만족할 기준)
% ====================================================================
\section*{부록 B --- 코드 구현 요구명세(회귀 기준$\cdot$하위호환)}
\addcontentsline{toc}{section}{부록 B --- 코드 구현 요구명세(회귀 기준·하위호환)}

★\textbf{doc-leads.} 이 문건이 권위이며, 코드(\code{Anode\_Fit})는 이후 이 문건에 맞춰 개정된다. 따라서 이
부록은 \emph{현 구현을 기술하지 않는다} --- 대신 코드 개정이 \emph{재현해야 할} 회귀 기준값과 하위호환 조건을
요구명세로 못박는다. 본문 물리$\cdot$식이 상위이고, 아래는 그 물리를 옮긴 구현이 통과해야 할 검사다. 함수명은
이 부록에만 등장한다(본문은 교과서 register 로 물리$\cdot$모델 언어만 쓴다).

\textbf{B.1 최종 구현 타깃.} 실전 계산의 유일 타깃은 종합식~\eqref{eq:complete} 이다: 입력
$\{\Delta S^0_j,\,Q_j,\,U_j,\,w_j\}$ 와 Chapter 1 의 $\xi_j(V,T)$~\eqref{eq:logistic} 를 받아, 전하 보존
음함수~\eqref{eq:implicit} 로 $U_\oc(\bar x,T)$ 를 풀고, $\partial U_\oc/\partial T$ 와 $\dot Q_\rev=-IT\,\partial
U_\oc/\partial T$~\eqref{eq:qrev} 를 낸다. 진동 온도 편차가 지정되면 입력을 $\Delta S^0_j\to\Delta
S^0_j+\Delta S_\vib(T)$, $U_j\to U_j+\Delta U_\vib(T)$ 로 온도 확장한다(\S\ref{sec:einstein}).

\textbf{B.2 회귀 기준값.} 아래 값을 구현이 재현해야 한다(298.15 K, Chapter 1 4-전이 staging, $w_j{=}RT/F$).
\begin{longtable}{p{0.30\textwidth} p{0.42\textwidth} c}
\toprule
요구 항목 & 재현해야 할 기준값 & 근거 \\
\midrule
\endfirsthead
\toprule
요구 항목 & 재현해야 할 기준값 & 근거 \\
\midrule
\endhead
\code{entropy\_coefficient}$(\bar x{=}0.25)$ &
$\partial U_\oc/\partial T=-0.204$ mV/K(완전식); 단순식 하한 $-0.134$; config 몫 $-0.070$ &
\S\ref{ssec:worked}(c) \\
$U_\oc(\bar x{=}0.25)$ &
$74.4$ mV(음함수~\eqref{eq:implicit} 유일근) &
\S\ref{ssec:worked}(a) \\
round-trip 유한차분 &
$[U_\oc(T{+}3){-}U_\oc(T{-}3)]/6\,\mathrm K=-0.204$ mV/K, 해석 완전식과 $<0.001\,\mu$V/K 일치 &
\S\ref{ssec:worked}(d) \\
$\dot Q_\rev/I(\bar x{=}0.25)$ &
$\Delta S=-19.7$ J\,mol$^{-1}$K$^{-1}$, $\dot Q_\rev/I=+60.8$ mV(방전 발열) &
\S\ref{ssec:worked}(e) \\
SOC 부호 교대(5점) &
$\bar x{=}0.10/0.25/0.50/0.75/0.90$ 에서 $\partial U_\oc/\partial T=-0.307/-0.204/-0.089/+0.044/+0.218$ mV/K,
$\dot Q_\rev/I=+91.5/+60.8/+26.6/-13.2/-64.9$ mV(발열$\to$흡열) &
표~\ref{tab:qrev} \\
전 범위 파생 A &
175 점에서 완전식이 유한차분과 표시 정밀도($\lesssim10^{-2}$ mV/K) 일치, 내부 경계 3 곳 연속 블렌드 &
\S\ref{ssec:blend}, \cite{numverif2026} \\
\bottomrule
\end{longtable}

\textbf{B.3 하위호환 조건.} Einstein 진동 보정(\S\ref{sec:einstein})은 \emph{선택적}이어야 한다: 전이 $j$ 의
Einstein 온도 $\theta_{E,j}$ 를 지정하지 않으면 두 가산 $\Delta U_\vib(T)$, $\Delta S_\vib(T)$ 가 모든 $T$ 에서
항등적으로 $0$ 이 되어, 보정을 켜기 이전 계산과 \emph{완전히 동일한 값}(bit-exact)을 내야 한다. 곧 이 개정은
$\theta_{E,j}$ 미지정 경로에서 회귀를 일으키지 않는다.

\textbf{B.4 입력 규약 준수.} 평형 $\xi_j(V,T)$ 는 방향 부호가 아니라 $s{=}+1$$\cdot$평형 중심 $U_j$(분기 shift
無)로 평가한다(\S\ref{ssec:synth}). 두-상 전이의 폭 $w_j$ 는 평형 예측 $n_jRT/F$ 를 강제 값으로 쓰지 않고
다온도 $dQ/dV$ 피팅의 자유 폭으로 받는다(\S\ref{ssec:weff}); 지배 두-상 config 몫은 다온도 round-trip 확정
전까지 단순식을 보수적 기준으로 둔다(\S\ref{ssec:synth}).

% 자산: (신규 물리 자산 없음 — doc-leads 요구명세 표면. 코드 anchor 재프레이밍:
%   [C-111] entropy_coefficient −0.204·round-trip 회귀 기준 · [C-114] tab:qrev 5점 회귀값 ·
%   [C-61] θ_E 미지정 bit-exact 하위호환 · [C-102/103] eq:complete 최종 타깃·입력 인터페이스 · [C-133] numverif 전범위 기준)
